\section{Optimized short comparisons}\label{sec:comparisons}

This section proves local estimates for a constrained comparison, including
an error estimate conditional on regularity outside the scale being tested.
The regularity assumption is stated each time it is used.  The later
propagation argument establishes it by a finite sequence of further
comparisons; none of the estimates below assumes the conclusion of that
argument.
The construction uses the optimized interpolation and pinned-field
averaging framework of \cite[Proposition~2.3 and Section~3]{AL2026}, with
the local strength profile and the scale-dependent constraints retained
throughout.

\subsection{Parameters and the constrained comparison}

Fix a sufficiently small constant $m>0$ and put $\delta=m/10000$.
All asymptotic statements concern $n\to\infty$ with these constants fixed.
The temperature, bounds on the total clocks, and fixed comparability
constants may enter implicit constants.  A polynomial saving means a
factor $n^{-c}$ for some fixed $c>0$.  Such a saving, unless otherwise
specified, is independent of each fixed moment order; the multiplicative
constant and the threshold in $n$ may depend on that order.  Factors denoted
by $n^{o(1)}$ can be bounded by $C_\epsilon n^\epsilon$ for any fixed
$\epsilon>0$.  We choose $\epsilon$ after the positive power margins in an
estimate have been fixed.  The estimates hold along arbitrary sequences
of parameters satisfying the stated assumptions, and hence also along
subsequences. We write $o_{\rm poly}(A)$ for a quantity bounded in
absolute value by $Cn^{-c}A$ for some fixed $c>0$. A bound $Y=O_p(A)$
means $\|Y\|_{L^p}\le C_pA$ at each fixed moment order under discussion.

A comparison call has parameters
\begin{equation}
 n^{-2/3+m}d\lesssim E\le n^{-1/2}d^2,
 \qquad a=\frac{\sqrt n E}{d},\qquad
 u_*=\frac1{nE},\qquad 0<d=O(1).
 \tag{3}\label{eq:3}
\end{equation}
The parameter $E$ sets the comparison-error budget, up to the fixed
power slack used below. The maximal covariance-removal strength is
$d^2$; at that strength the profile in \eqref{eq:6} changes scale near
radius $d$. The quantity $a$ is the corresponding lower direct-test
scale, while $u_*$ controls the smallest-mass constraints.
In particular $d\ge a\gtrsim n^{-1/6+m}$.  The working interval of masses
is $[u_0,1]$, where either $u_0=0$ or $u_0=o(a)$; changes made by this call
vanish below $u_0$.  Its endpoint clocks are denoted by $b_0,h_0$.  They are
uniformly bounded and nondecreasing on the working interval.  For the
subsequent profile argument the endpoint hypothesis is
\begin{equation}
 b_0(u)=T+o(z^2),\qquad h_0(u)=o(z^3)
 \quad\text{when }u\asymp z,\quad a\lesssim z\le1.
 \tag{4}\label{eq:4}
\end{equation}
This formulation is sequential: it holds on every such test sequence,
with any fixed comparability constants and with masses in the open
working interval.  In particular it implies uniform little-oh estimates
on each fixed compact scaled mass band, by choosing a violating point if
uniformity were to fail.

Choose $r_{\min}=0$ in the all-mass case.  In the case $u_0>0$ we may fix
any $0\le r_{\min}\le u_*/M(u_0)$.  The optimization variable is an
increasing onto quantile $r:[u_0,1]\to[r_{\min},1]$.  Its inverse
$\alpha$ satisfies $\alpha(r_{\min})=u_0$, $\alpha(1)=1$, and
\begin{equation}
 \begin{gathered}
 k(r)\le\alpha'(r)\le M(\alpha(r)),\qquad
 k(r)=\frac{n^\delta u_*^2}{(r+n^{2\delta}u_*)^2},\\
 M(u)=n^\delta\left[\frac{u+u_*}{L}
                  +\frac{d^3}{E}(u+u_*)^2\right],
 \qquad L=\min\left(u_*,\frac{a^3}{d^2}\right).
 \end{gathered}
 \tag{5}\label{eq:5}
\end{equation}
These are almost-everywhere constraints on Lipschitz functions, with a
strictly positive floor.  The inverse is therefore Lipschitz at fixed
$n$ as well.  The parameters in \eqref{eq:5}, including $r_{\min}$, remain
fixed when the strength is varied.

Fix once and for all a smooth function $f>0$ that equals $1$ in a
neighborhood of $0$, is comparable to $(1+\log t)/t$ for large $t$, and
whose logarithmic slope is nondecreasing from $0$ to $1$.  We require
$1-p\gtrsim[1+\log(1+t)]^{-1}$ for that slope.  Such a function is
obtained by choosing a smooth nondecreasing slope that is zero near zero
and equals $1-(1+\log t)^{-1}$ for all sufficiently large $t$, and then
integrating its logarithmic derivative.  Set
\begin{equation}
 \begin{split}
 e(r)&=\lambda f(\lambda r/d^3),\qquad 0\le\lambda\le d^2,\\
 \mathcal F_\lambda(r)
 &=f_n\bigl(b_0-e(r),h_0+e(r)r\bigr)
       +\frac14\int_{u_0}^1 e(r(u))r(u)^2\,du,
 \qquad \mathcal M_\lambda=\min_r\mathcal F_\lambda(r).
 \end{split}
 \tag{6}\label{eq:6}
\end{equation}
Here and below $e(r)$ in a clock denotes composition with $r(u)$, and
the changes in \eqref{eq:6} are zero below $u_0$.  We write
$p(r)=-re'(r)/e(r)$.  Thus $0\le p<1$,
$\partial_{\log\lambda}e=(1-p)e$, and $(re)'=(1-p)e\ge0$.
Consequently the comparison preserves monotonicity inside the working
interval.  A call at a fixed strength is admissible when the changed
clocks in \eqref{eq:6} are nonnegative and nondecreasing for every
quantile in \eqref{eq:5}, including across the boundary $u_0$.
An admissible sweep requires this for every strength in the sweep,
including zero.  The endpoints used later will satisfy these conditions.
A bare endpoint may instead have a field discontinuity at $u_0$ that
is repaired only at a fixed positive strength.  In that case the local
estimates concern the changed paths at that strength; the bare pressure
and a sweep from zero are not used.

\begin{lemma}[Existence and the optimized envelope]\label{comp:envelope}
For all sufficiently large $n$, the class \eqref{eq:5} is nonempty and
compact in the uniform topology of inverse quantiles.  The minimum in
\eqref{eq:6} is attained.  For an admissible sweep, $\mathcal M_\lambda$
is locally Lipschitz and nondecreasing in $\lambda$, equals the endpoint
pressure at $\lambda=0$, and is an upper bound for that pressure.
At almost every positive strength its upper derivative per unit
$\log\lambda$ is bounded, at any minimizer at that strength, by
\begin{equation}
 \frac14\int_{u_0}^1 e(r(u))
           \bigl[D_u^2+(S_u-r(u))^2\bigr]\,du.
 \tag{7}\label{eq:7}
\end{equation}
All moments in this formula are for the changed clocks in
\eqref{eq:6}.
\end{lemma}

\begin{proof}
We have $k\le n^{-3\delta}$ and $M\ge n^\delta$, while
\[
 \int_0^1 k(r)\,dr\le n^{-\delta}u_* =o(1),\qquad
 \int_{u_0}^1\frac{du}{M(u)}
 \le n^{-\delta}L\log\frac{1+u_*}{u_0+u_*}=o(1).
\]
The floor alone does not reach the prescribed terminal mass.  The cap
flow reaches it before the prescribed terminal radius.  Interpolating
controls in $\alpha'=k+w(M(\alpha)-k)$, $0\le w\le1$, gives a feasible
endpoint by continuity of the scalar flow; equivalently, interpolate the
travel time until the cap flow reaches mass $1$.  At fixed $n$ the
derivatives are uniformly bounded by $\max_{[u_0,1]}M$.  Uniform limits
preserve both differential inequalities in their integrated form, so
Arzel\`a--Ascoli gives compactness.  The positive floor preserves
invertibility.  Continuity of the pressure follows from the bounded
variation formula \eqref{eq:1}; this proves attainment.

For a fixed quantile, interpolate the amount of covariance removed from
$0$ to $e$.  The derivative of the pressure plus the corresponding
penalty is the integral of $[B-2rS+r^2]/4$ times the nonnegative strength
increment.  Since $B=S^2+D^2$, it is nonnegative.  This proves the
upper-bound assertion and monotonicity, including after minimizing.
At positive strength use an old minimizer as a competitor for a positive
increment of strength.  The same differentiation and
$\partial_{\log\lambda}e=(1-p)e\le e$ give \eqref{eq:7}.
Bounded spin moments give local Lipschitz continuity of the objective
uniformly over the compact admissible class, and hence of its minimum.
Integration of its almost-everywhere derivative is therefore legitimate.
The whole interval $0\le\lambda\le n^{-10}$ costs $O(n^{-10})$.
\end{proof}

\begin{definition}[Outer regularity]\label{comp:outer-regularity}
A minimizer has a regular outer region above $x$ if, for fixed positive
constants $c,C,C_0$, one has $cu\le r(u)\le Cu$ whenever
$u\ge C_0x$.  Enlarging the cutoff by a fixed factor gives the
equivalent comparison $\alpha(r)\asymp r$ in radius coordinates.  The
constants are independent of $n$.  No bound on a pointwise path
derivative is part of this definition.
\end{definition}

The eventual regularity assertion concerns every test scale
\begin{equation}
 n^\delta\frac{\sqrt n E}{\sqrt\lambda}\le x\le1.
 \tag{8}\label{eq:8}
\end{equation}
It will be proved after the profile analysis, by propagation from larger
scales.  In the present section it is a target range, not an assumption
already established at every minimizer.  Write
$A_\lambda=n^\delta\sqrt n E/\sqrt\lambda$.
If $x\ge A_\lambda$, then $x\gtrsim a$ and
$e(x)x^2\gtrsim n^\delta nE^2$: indeed
$\sqrt\lambda\sqrt n E\le d^3$, so the argument of $f$ at
$A_\lambda$ is at most $n^\delta$, and $re(r)$ increases.
Our local estimates require only $x\gtrsim a$ and
$e(x)x^2\gtrsim nE^2$, together with outer regularity when $x=o(1)$.

\subsection{First-variation geometry}

Define
\[
 G(r)=e(r)\left[(2-p(r))r-2(1-p(r))S_r-p(r)B_r/r\right],
\]
with the final term interpreted as zero in the neighborhood where $p=0$.
Subscripts $r$ here mean revelation at mass $\alpha(r)$.

\begin{lemma}[Contact, cap, and floor]\label{comp:geometry}
At a minimizer there is a continuous absolutely continuous potential
$H$ such that $H'=JG$ for a strictly positive factor $J$,
$\alpha'=M(\alpha)$ on $\{H<0\}$, and $\alpha'=k$ on $\{H>0\}$.
On a band where $\alpha+u_*$ varies by a bounded factor, $J$ has bounded
distortion after normalization at one endpoint.  On the contact set
$\{H=0\}$, almost everywhere,
\begin{equation}
 G=0,\qquad S'\lesssim n^{o(1)},\qquad S\le r.
 \tag{9}\label{eq:9}
\end{equation}
On an open cap component $\{H<0\}$ of radius length $\ell$,
\begin{equation}
 \begin{gathered}
 \Delta S\lesssim n^{o(1)}
        (\ell+r_{\min}\mathbf1_{\mathrm{initial}}),
 \qquad S\le r+\ell,\\
 |S-r|^2\lesssim n^{o(1)}
   \left[D^2+\ell^2+r_{\min}^2\mathbf1_{\mathrm{initial}}\right].
 \end{gathered}
 \tag{10}\label{eq:10}
\end{equation}
Here $\Delta S$ is the increase through the interior of the component.
On contact, $(S-r)^2\lesssim D^2$.  Every cap has
$\ell\lesssim n^{-\delta}L\log n$; if it is contained in a comparable
mass band $u\asymp v\gg u_*$, then
$v\ell\lesssim n^{-\delta}E/d^3$.  A noninitial cap on which $p=0$
throughout satisfies $|S-r|\le\ell$.  The total floor mass is at most
$n^{-\delta}u_*$.
\end{lemma}

\begin{proof}
Differentiating the objective with respect to its forward quantile gives
$\frac14\int G\,\Delta r\,du$.  The inverse variation is therefore
$-\frac14\int G\,\Delta\alpha\,dr$.  Write
$\alpha'=k+w(M(\alpha)-k)$ with $0\le w\le1$.  Its linearized
equation has fundamental factor
\[
 J(r)=\exp\left(\int_{r_{\min}}^r w(t)M'(\alpha(t))\,dt\right).
\]
Terminal preservation says
$\int(M-k)\Delta w/J=0$.  Integrating the first variation in the
opposite order gives a box-constrained linear functional with this one
linear constraint.  A scalar multiplier gives a function with derivative
$JG$, which must be nonpositive where $w=1$ and nonnegative where $w=0$.
Both signs of correcting endpoint variation exist: a feasible path
cannot be the pure floor or pure cap, as their travel times above show.
Thus the multiplier can equally be obtained by adding a correcting
variation and taking one-sided limits.  This proves the asserted rule.
Directional derivatives follow by the segment variation formula; the
fixed-$n$ inverse functions have dominated difference quotients.  Values
at jumps affect no integrals.  Moreover
\[
 \frac{M'(u)}{M(u)}\le\frac2{u+u_*},\qquad wM\le\alpha',
\]
which proves bounded distortion of $J$ on the stated bands.

The bracket in $G$ decreases with $p$, since its $p$ derivative is
$-[(r-S)^2+D^2]/r$.  At a noninitial left cap endpoint its limiting
value is nonpositive, and at the right endpoint it is nonnegative.
At radius $1$ the latter remains true because $S,B\le1$.
The right inequality implies $S\le r$ there and $pB\le2r^2$.
Use the right-end slope parameter at both endpoints.  If their radii
are comparable, subtraction of the inequalities, monotonicity of $B$,
and $pB\le2r_{\mathrm{right}}^2$ give
$2(1-p)\Delta S\le C\ell$.
If the radii are not comparable, $S\le r_{\mathrm{right}}\lesssim\ell$
suffices.  At an initial component use $S\ge0$, introducing only
$r_{\min}$.  Since $(1-p)^{-1}=O(\log n)$ on the polynomial scales
under consideration, this proves the first line of \eqref{eq:10}.

For the mismatch bound, the case $S>r$ follows from the right endpoint.
In a noninitial component with comparable endpoint radii, the left
inequality and monotonicity, followed by replacement of its radius by
the current radius, give
\[
 (r-S)\left[(2-p)+pS/r\right]\le C\ell+pD^2/r.
\]
When $S<r$, this gives $r-S\le C(\ell+D^2/r)$; splitting into
$D\ge r$ and $D<r$ proves the squared estimate.  Noncomparable and
initial components are covered by their length and wall terms.
The same calculation at $G=0$ proves $S\le r$ and
$(S-r)^2\lesssim D^2$ on contact.  At density points of the zero set,
difference quotients between adjacent zero-set points yield
$S'\lesssim(1-p)^{-1}$: the derivatives of $B$ and $p$ have favorable
signs, and $pB/r^2\le2$ there.  This proves \eqref{eq:9}.
The exceptional contact set has zero radius length and, since $\alpha$
is Lipschitz at fixed $n$, zero mass.

On a cap, $dr=du/M(u)$, so the two length bounds follow respectively
from the two positive terms of $M$.  If $p=0$, the endpoint signs say
$S_{\mathrm{left}}\ge r_{\mathrm{left}}$ and
$S_{\mathrm{right}}\le r_{\mathrm{right}}$, proving the sharper
mismatch statement.  Finally, floor mass is bounded by
$\int_0^1 k(r)\,dr\le n^{-\delta}u_*$.
\end{proof}

\subsection{Projection estimates and the raw-moment bootstrap}

Assume for now that $x=o(1)$, $x\gtrsim a$,
$e(x)x^2\gtrsim nE^2$, and the minimizer has a regular outer region
above $x$.  After a fixed enlargement of the cutoff, all bands below
are regular in both coordinates.  On the working interval
$\partial_r[h(\alpha(r))]\ge e(r)/O(\log n)$.
Hereafter $\partial_r h$ denotes this derivative of the clock pulled
back to radius coordinates. At regular scales $r\asymp u\asymp v$,
\begin{equation}
 u_*\ll x,\qquad
 M(O(v))\ll(nEv)^2\lesssim n e(v)v^4,
 \qquad n^{-1/6}\ll x.
 \tag{11}\label{eq:11}
\end{equation}
Every strict gain in this display is polynomial.  Indeed
\[
 \frac{u_*}{a}\lesssim n^{-3m},\qquad
 \frac{u_*^2d^2}{a^4}\lesssim n^{-6m},\qquad
 \frac{d^3}{n^2E^3}\lesssim n^{-3m}.
\]
For $v\gtrsim a$, dividing $M(O(v))$ by $(nEv)^2$ bounds the ratio by
$Cn^\delta$ times the sum of these three quantities.
The comparison with $ne(v)v^4$ follows because $e(r)r^2$ increases.
Monotonicity and the outer bounds also show that if one of $u,r$ is
$O(x)$, the other is $O(x)$.

\begin{proposition}[Raw moments]\label{comp:raw-moments}
For every fixed $p\ge2$, throughout the path, including masses below
the working interval,
\begin{equation}
 \|X_u\|_p\lesssim_p\sqrt{u+x},\qquad
 \|R_{\mathrm{split\ at}\ u}\|_p\lesssim_p u+x.
 \tag{12}\label{eq:12}
\end{equation}
The bounds also hold at intermediate clock revelations, with their
specified mass.  The norm of a random vector is its Euclidean norm
inside the $L^p$ norm.
\end{proposition}

We first establish a projection estimate that will also be used later.

\begin{lemma}[Projection after a field-clock interval]\label{comp:projection}
Let $V$ be a vector determined by outside branches and by data at or
before their fork with a branch $b$.  Conditional on this data, the
future of $b$ has its usual planted law.  Suppose a field-clock interval
after the fork has length $D_h$, with mass at least $v_->0$ throughout.
For even fixed $p\ge2$, the part of $V\cdot y_b$ left after projection
at the end of the interval satisfies
\begin{equation}
 \|V\cdot(y_b-X_{b,\mathrm{end}})\|_p
 \le C_p(\log n)^2\frac{\|V\|_p}{v_-\sqrt{nD_h}}.
 \tag{13}\label{eq:13}
\end{equation}
Here the masses and clock lengths used in the estimate are bounded
below by inverse powers of $n$.  Differences of two later projections
satisfy the same estimate up to a constant.
\end{lemma}

\begin{proof}
Condition on the outside data so that $V$ is fixed for the branch
calculation.  Put $Y=V\cdot y_b$ and $Y_t=V\cdot X_{b,t}$.
In field-clock time, the spin Brownian driver has planted drift
$t\sqrt n X_{b,t}$.  This follows on each filled increment from the
normalized exponential tilt and its logarithmic gradient; subdivision
gives the Brownian formulation.  Gaussian differentiation of the
telescoping planted density gives, for a subinterval $I$ and
$P=(Y-Y_{\mathrm{start}})^{p-1}$,
\[
 \mathbb E\left[P\int_I V\cdot(dG-t\sqrt n X_{b,t}\,dh(t))\right]
 =\sqrt n\int_I\mathbb E\left[
 P\left\{t(Y-Y_t)+\int_t^1(Y-Y_z)\,dz\right\}\right]dh(t).
\]
The terminal tilt contributes $Y$; the subsequent continuation terms
and the subtracted innovation drift contribute
$-\int_t^1Y_z\,dz-tY_t$.  Prefix-dependent parts of the test are fixed
before the subinterval, so no derivative of $V$ or
$Y_{\mathrm{start}}$ is taken.  This identity is first an ordinary
finite Gaussian integration by parts and then passes to filled clocks
by the path approximation already established.

For $z$ after the subinterval's start, conditional centering gives
\[
 \mathbb E[P(Y-Y_z)]
 =\mathbb E\left[
 \bigl((Y-Y_{\mathrm{start}})^{p-1}
       -(Y_z-Y_{\mathrm{start}})^{p-1}\bigr)(Y-Y_z)\right]
 \ge c_p\mathbb E|Y-Y_z|^p.
\]
This is the uniform monotonicity inequality for the odd power $p-1$.
Also $\|Y-Y_{\mathrm{end}}\|_p\le2\|Y-Y_t\|_p$ for earlier $t$.
The innovation on the left, divided by $\sqrt{h(I)}$, has norm at
most $C_p\|V\|_p$.  H\"older's inequality therefore gives the recursion
\[
 \|Y-Y_{\mathrm{end}}\|_p^p
 \le \frac{C_p\|V\|_p}{v_-\sqrt{nh(I)}}
       \|Y-Y_{\mathrm{start}}\|_p^{p-1}.
\]
Divide the given clock into $N=\lceil(\log n)^2\rceil$ equal pieces.
The coefficient is
\[
 \frac{C_p\sqrt N\|V\|_p}{v_-\sqrt{nD_h}}.
\]
Iteration raises the ratio of the initial bound $2\|V\|_p$ to this
coefficient to the power $(1-1/p)^N$.  If the ratio is larger than one
it is at most a fixed power of $n$, so its final power tends to one.
If it is smaller than one the assertion is already immediate.
This proves \eqref{eq:13}, with its harmless logarithmic allowance.
Conditional Jensen proves the assertion for later projections.
\end{proof}

\begin{proof}[Proof of Proposition~\ref{comp:raw-moments}]
\emph{Integrated susceptibility bound.}
Fix an integer $l\ge1$ and let
\[
 Q=\max\left(1,\sup_u
       \frac{\|X_u\|_{4l}}{\sqrt{x+u}}\right),
\]
including intermediate revelations and the final spin.  This is finite
at each fixed $n$.  We first bound the integrated weight in
\eqref{eq:1} on a small regular band of mass and radius comparable to
$v$.  Uniformly on a bounded enlargement of this band, choose $\rho$ so
that
\begin{equation}
 M(O(v))\rho\le n^{-c},\qquad
 ne(v)v^4\rho\ge n^c
 \tag{14}\label{eq:14}
\end{equation}
for a fixed $c>0$ independent of $l$.  For example take the geometric
mean between the reciprocal scales; \eqref{eq:11} gives the required
room, and $\rho=o(1)$.

We give the layer estimate first on a finite filled hierarchy, appending
the terminal spin as a martingale step of mass $1$.  In this paragraph
parentheses denote radius coordinates: $X(s)=X_{\alpha(s)}$,
$\mathcal F(s)=\mathcal F_{\alpha(s)}$, and likewise $H(s),W(s)$;
write $X(\dagger)=y$ and $\mathbb E_s=\mathbb E[\,\cdot\mid\mathcal F(s)]$.
Enclose the band in a radius interval $I=[r_-,r_+]$, with
$r_\pm\asymp v$ and $r_+<r_c/2$ for a fixed small $r_c>0$.
Its fixed enlargement is regular.  Put
\[
 \mathsf A_v=M(Cv),\quad \mathsf B_v=ne(v)v^4,\quad
 \rho=(\mathsf A_v\mathsf B_v)^{-1/2},\quad \eta=\rho v.
\]
By \eqref{eq:11}, with $\kappa=3m-\delta>0$,
$\mathsf A_v/\mathsf B_v\lesssim n^{-\kappa}$; hence
$\mathsf A_v\rho=(\mathsf B_v\rho)^{-1}\lesssim n^{-\kappa/2}$, as required
in \eqref{eq:14}.  Also $\mathsf A_v\ge n^\delta$, so $\rho=o(1)$.

For each $s\in I$, fork an independent continuation $b$ from the original
branch $a$.  On both branches use cuts
$t_0=s$, $t_1=s+\eta$, $t_2=2s$, then doubled radii up to the first
$t_J\in[r_c,2r_c)$, and $t_{J+1}=\dagger$.
Set $\Delta_i^a=X^a(t_{i+1})-X^a(t_i)$, and similarly for $b$.
There are $K=J+1=O(\log n)$ layers.  With
$w_0=w_1=s$, $w_i=t_i$ for $2\le i<J$, and $w_J=1$,
their largest masses $\mu_i$ satisfy
$\mu_i\lesssim w_i$ and $\|\Delta_i^{a,b}\|_{4l}\lesssim_l Q\sqrt{w_i}$.
For elementary martingale increments $d_\nu^b$ after the fork,
conditional orthogonality gives the indexed positive-semidefinite bound
\[
 H(s)=\mathbb E_s\sum_{\nu\text{ after }s}
        m_\nu d_\nu^b(d_\nu^b)^{\mathsf T}
 \preceq\sum_{j=0}^J\mu_j
        \mathbb E_s[\Delta_j^b(\Delta_j^b)^{\mathsf T}].
\]
Indeed cross terms between distinct elementary increments in each
block have zero conditional expectation.  Since
$y^a-X(s)=\sum_i\Delta_i^a$, scalar Cauchy--Schwarz yields
\[
 W(s)\le K\sum_{i,j=0}^J\mu_j\,
   \mathbb E_b[(\Delta_i^a\cdot\Delta_j^b)^2\mid a].
\]
Here conditioning on $a$ includes the common prefix; no coupling of the
auxiliary branches for different $s$ is needed.

For $(i,j)\ne(0,0)$ choose the branch carrying layer
$k=\max(i,j)$, taking either branch on a tie.  Condition on the entire
opposite branch and use its increment as the exterior anchor in
Lemma~\ref{comp:projection}.  The clock preceding the selected layer is
\[
 [s,s+\eta]\quad(k=1),\qquad
 [3t_k/4,t_k]\quad(2\le k\le J).
\]
Its masses are bounded below by $cw_k$, also for the terminal layer
because $t_J\asymp1$.  The field-clock length is at least
$c\rho ve(v)/\log n$: for $k=1$ this follows from
$\partial_r h\gtrsim e(r)/\log n$, and for $k\ge2$ from the same bound and
monotonicity of $re(r)$, with no factor $\rho$ needed.
All these scales are bounded below by inverse powers of $n$.
The opposite increment has norm at most $C_lQ\sqrt{w_k}$, even when
its endpoint is later than the fork, since it belongs to an independent
branch.  Equation~\eqref{eq:13} and conditional Jensen for later
projections therefore give
\[
 \mu_j\|\Delta_i^a\cdot\Delta_j^b\|_{2l}^2
 \le \frac{C_l(\log n)^5Q^2}{\rho nve(v)}.
\]
Conditional Jensen also bounds the $L^l$ norm of the corresponding
conditional square by this squared $L^{2l}$ norm.  Integrating over
physical mass $O(v)$ and summing the $K^2$ pairs, including the factor
$K$ above, bounds their total by
$C_l(\log n)^8Q^2v^4/(\mathsf B_v\rho)$.

For $(i,j)=(0,0)$ put
$q_s=\mathbb E_s|X(s+\eta)-X(s)|^2$ and
$Z=\sup_{r_-\le t\le r_++\eta}|X(t)|$ on the original path.
The integrand is at most $CvZ^2q_s$; changing coordinates costs only
$du\le \mathsf A_v\,ds$.  Doob's inequality gives
$\|Z^2\|_{2l}\le C_lQ^2v$.  To estimate the remaining integral, for
$\xi\in[0,\eta)$ let $r_j=r_-+\xi+j\eta$, retaining precisely the
$N=N(\xi)$ starting points in $I$ and their last endpoint
$r_N\le r_++\eta$.  Set
\[
 U_j=|X(r_j)|^2,\quad A_0=0,\quad
 A_{j+1}=A_j+q_{r_j},\qquad
 q_{r_j}=\mathbb E[U_{j+1}-U_j\mid\mathcal F(r_j)].
\]
Thus $A_{j+1}$ is $\mathcal F(r_j)$-measurable.  For $p=2l$,
$A_N^p\le p\sum_jA_{j+1}^{p-1}q_{r_j}$, and the exact identity
\[
 \begin{split}
 \sum_{j=0}^{N-1}A_{j+1}^{p-1}(U_{j+1}-U_j)
 ={}&A_N^{p-1}U_N-A_1^{p-1}U_0\\
 &-\sum_{j=1}^{N-1}(A_{j+1}^{p-1}-A_j^{p-1})U_j\\
 \le{}&A_N^{p-1}U_N
 \end{split}
\]
gives $\mathbb E A_N^p\le p\mathbb E[A_N^{p-1}U_N]$.
Consequently $\|A_N\|_{2l}\le2l\|U_N\|_{2l}\le C_lQ^2v$,
independently of the grid cardinality.  The shifted-grid identity
\[
 \int_Iq_s\,ds
   =\int_0^\eta\sum_{j=0}^{N(\xi)-1}q_{r_-+\xi+j\eta}\,d\xi
\]
and Minkowski give $\|\int_Iq_s\,ds\|_{2l}\le C_l\eta Q^2v$.
H\"older therefore bounds the integrated first-pair term by
$C_lv\mathsf A_v(Q^2v)(\eta Q^2v)=C_l\rho \mathsf A_vQ^4v^4$.

Combining the pairs, including their logarithmic multiplicities, gives
\[
 \left\|\int_{\mathrm{band}}W_u\,du\right\|_l
 \le C_l(\log n)^8
       \left[\mathsf A_v\rho Q^4+(\mathsf B_v\rho)^{-1}Q^2\right]v^4.
\]
All estimates use moments at most $4l$.  Include deterministic jumps
and intermediate cuts in the finite hierarchy.  An after-revelation
endpoint assigns its jump to the preceding block; a before-revelation
endpoint assigns it to the following block.  Equivalently insert both
ordered nodes, counting the jump exactly once.  The constants are independent of the
mesh and grid cardinality, so the established bounded-path approximation
passes to filled clocks at fixed $n$; the countable set of jump radii
does not affect the radius integral.  Absorbing logarithms, for example
with $c'=\kappa/4$, proves
\begin{equation}
 \left\|\int_{\mathrm{band}}W_u\,du\right\|_l
 \lesssim_l n^{-c'}(Q^4+Q^2)v^4,
 \tag{15}\label{eq:15}
\end{equation}
where $c'>0$ is independent of $l$.  The argument also bounds integrals
over subsets, since $W\ge0$.

\emph{Closing the moment estimate.}
For a regular small mass $u$, choose a future window of mass comparable
to $u$ at radii comparable to $u$.  Its means satisfy
$\mathbb EL_j=S_j/2\le Cu$.  Indeed this follows on nonfloor points
from Lemma~\ref{comp:geometry}; the floor has total mass $o(u)$, so
monotonicity permits looking ahead to a nonfloor point.  For $j\ge u$,
$2\mathbb E_uL_j=\mathbb E_u|X_j|^2\ge|X_u|^2$.
Apply \eqref{eq:1} and \eqref{eq:15} to the normalized window average.
Conditional Jensen yields
\[
 \||X_u|^2\|_{2l}\le Cu+C_l n^{-c''}Q^2u.
\]
For smaller masses, including $u<u_0$, choose a deterministic future
regular cut $t\asymp x$.  The identity
$X_u=\mathbb E[X_t\mid\mathcal F_u]$ and conditional Jensen give
\[
 \|X_u\|_{4l}^2\le\|X_t\|_{4l}^2
                    \le Cx+C_l n^{-c''}Q^2x.
\]
Macroscopic masses are bounded directly by $|X|\le1$.
Taking the supremum gives $Q^2\le C_l+C_l n^{-c''}Q^2$, which closes
for all sufficiently large $n$.  This proves the first bound of
\eqref{eq:12} for exponents divisible by four, and hence for every
fixed exponent.  It also removes $Q$ from \eqref{eq:15}.

\emph{The two-leaf overlap.}
If $u+x$ is bounded below, the overlap bound is immediate from $|R|\le1$.
Otherwise choose a regular radius $q=C(u+x)$, with $C$ fixed and
sufficiently large that $[q/2,q]$ is regular and lies after the fork.
Outer comparability and monotonicity permit this choice even when
$u<u_0$.  Write $X^a(q)=X^a_{\alpha(q)}$ and similarly for branch $b$.
The projected inner product has norm
$\|X^a(q)\cdot X^b(q)\|_p=O_p(q)$ by the bound just proved.

Decompose each leaf into its projection at $q$ and subsequent doubling
layers, ending with the terminal spin as in the susceptibility estimate.
Every product other than the two initial projections contains a later
layer of scale $w\ge q$.  An innovation interval before that layer has
mass at least $cw$ and field-clock length at least
$ce(w)w/\log n$.  The opposite branch supplies an exterior anchor of
norm $O_p(\sqrt w)$.  Lemma~\ref{comp:projection}, now without a
mass weight, bounds the product by
$C_p(\log n)^C/\sqrt{ne(w)w^2}$.
Since $e(w)w^2$ increases, summing the logarithmically many layer pairs
gives
\[
 \|R-X^a(q)\cdot X^b(q)\|_p
 \le\frac{C_p(\log n)^C}{\sqrt{ne(q)q^2}}
 =q\,\frac{C_p(\log n)^C}{\sqrt{ne(q)q^4}}
 =o_{\rm poly}(q),
\]
where \eqref{eq:11} supplies the last saving.  This proves the second
bound in \eqref{eq:12}.  All innovation intervals lie in the regular
region, including when the original fork is below the working interval.
\end{proof}

\begin{lemma}[Postprojection changes]\label{comp:resampling}
Suppose two leaves share only through a mass at most a fixed multiple
of $x_1$, where $x_1\ge x$.  If a projection or resampling is made after
a regular radius $v\ge Cx_1$, with comparable intervals fitting after
the fork and before that projection, the change in their inner product
caused by changing one endpoint has norm
\begin{equation}
 \frac{C_p(\log n)^C}{\sqrt{ne(v)v^2}}
 \tag{16}\label{eq:16}
\end{equation}
in every fixed $L^p$.  For an already projected anchor of norm
$O_p(\sqrt{x_1})$, the bound gains a factor $\sqrt{x_1/v}$.
\end{lemma}

\begin{proof}
Use the synchronized two-branch layer decomposition in the last part
of the preceding proof.  Every term involving the changed endpoint
contains an increment after its projection.  At the later layer of
scale $w$, the projection estimate gives
$C_p(\log n)^C/\sqrt{ne(w)w^2}$.  Monotonicity of $we(w)$ controls
the sum from its first scale $v$, up to logarithms.  For a fixed
projected anchor replace its raw $\sqrt w$ bound in
\eqref{eq:13} by $\sqrt{x_1}$.  Conditional Jensen covers replacement
by conditional projection, and two such replacements cover independent
resampling.
\end{proof}

\subsection{A split with uniform bounds in every fixed moment}

\begin{lemma}[Good split]\label{comp:good-split}
Suppose a regular band of mass and radius of order $v\ge Cx$ contains
a mass interval of length comparable to $v$ on which $S\asymp v$.
There are constants $c_1,c_2>0$, independent of each fixed moment order,
and a deterministic window of mass length $n^{-c_1}v$ with the following
properties.  Its fixed-factor enlargement has $S$ oscillation at most
$n^{-c_1/2}v$, and for every deterministic split in a fixed central
fraction of the window,
\begin{equation}
 \|R_*-S_*\|_p\le C_p n^{-c_2}v
 \tag{17}\label{eq:17}
\end{equation}
for all fixed $p\ge2$.  The window may be selected before the moment
orders or any later fixed recursion depth are chosen.  The constant
$c_1$ may be any sufficiently small positive constant.
\end{lemma}

\begin{proof}
Write $\theta=n^{-c_1}$, $\eta=n^{-c_1/2}$.  Among order
$\theta^{-1}$ disjoint buffered windows of mass length comparable to
$\theta v$, at most order $\eta^{-1}$ have $S$ oscillation larger
than $\eta v$, because the total increase is $O(v)$.  Choose one of
the remaining windows.  This is a deterministic selection involving
only $S$.  Fix past and future buffers of length comparable to
$\theta v$ surrounding a central region.

After the raw bootstrap, \eqref{eq:15} takes the form
$\|\int W\,du\|_q\le C_q n^{-\gamma}v^4$ for a $\gamma>0$
independent of fixed $q$.  Positivity of $W$ and \eqref{eq:1} give,
for the normalized $L$ average $A_J$ over either buffer,
\[
 \|A_J-\mathbb EA_J\|_p
 \le C_p n^{-\gamma/2+c_1}v=:C_p\zeta v.
\]
Let $t$ be a deterministic cut in the central region.  The exact
conditional identities give
\[
 2\mathbb E_t A_-\le|X_t|^2\le2\mathbb E_t A_+,
 \qquad
 |X_t|^2-2\mathbb E_t A_-
   =\operatorname{avg}_{s\in J_-}|X_t-X_s|^2\ge0.
\]
Since $2\mathbb EA_\pm$ differs from $S_t$ by at most $\eta v$,
conditional contraction implies
\[
 \||X_t|^2-S_t\|_p+
 \left\|\operatorname{avg}_{s\in J_-}|X_t-X_s|^2\right\|_p
 \le C_p(\eta+\zeta)v.
\]
For two central cuts $t,t'$ on the same branch, average the
triangle-square inequality through the common past buffer to obtain
\[
 \|X_t-X_{t'}\|_{2p}
 \le C_p(\eta+\zeta)^{1/2}\sqrt v.
\]
These estimates hold uniformly over deterministic cuts, including cuts
inside filled increments.  No supremum over random cuts is asserted.
Let $a_0=\min(c_1/2,\gamma/2-c_1)>0$.

Take the allowed star cuts in a smaller central fraction and project
the two branches at a later central cut separated by a fixed fraction
of $\theta v$.  This leaves a field clock at least
$c\theta e(v)v/[M(O(v))\log n]$ after every allowed fork.
For an outside anchor of norm $O_p(\sqrt v)$, \eqref{eq:13} bounds
the residual by
\[
 C_p(\log n)^C\frac{\sqrt{M(O(v))}}{\sqrt{n\theta e(v)}\,v}
 = C_p(\log n)^C v
       \sqrt{\frac{M(O(v))}{n\theta e(v)v^4}}.
\]
Decompose both leaves into the central projection and later layers.
The first later layers still have norm $O_p(\sqrt v)$, so the displayed
short clock applies.  For subsequent dyads use
Lemma~\ref{comp:resampling}; the later scales only improve its bound.
Thus the same bound, up to logarithms, controls the difference between
$R_*$ and the inner product of the two central projections.  This
two-branch decomposition is necessary: using a terminal spin of norm
one as the first anchor would lose a factor $\sqrt v$.

By \eqref{eq:11}, $M(O(v))/(ne(v)v^4)\le Cn^{-\kappa}$ for some
$\kappa>0$ independent of the moment order.  The preceding projection
distance estimates and H\"older give an error at most
$C_pn^{-a_0/2}v$ between the two central projections' inner product
and $|X_*|^2$.  Their square-norm concentration supplies the remaining
centering.  Choose $c_1$ sufficiently small.  For example any sufficiently
small
$c_2<\frac14\min(c_1,\gamma-2c_1,\kappa-c_1)$ works after absorbing
logarithms.  If extra fixed buffers cost a fixed power of $\theta$,
replace the last $c_1$ by that fixed multiple when making the choice.
All choices are independent of fixed $p$.
\end{proof}

\begin{remark}[Fixed-degree tree tests]\label{comp:good-split-trees}
The estimate \eqref{eq:17} needs no additional simultaneous-concentration
version when used in a fixed-degree tree test.  In every positive
planted law for a deterministic enlarged tree, pruning all other
branches leaves the ordinary two-leaf marginal of an edge splitting at
$*$.  A projection factor is a conditional average of such an edge.
Apply \eqref{eq:17} to these marginals and conditional Jensen to
projection factors, then use H\"older.  The finite-label counting rule
\eqref{eq:2} gives total variation bounded by a constant depending only
on the number of added labels, not on a minimum mass or the mesh.
Consequently a signed integral with $k$ new centered star factors and
a bounded remaining insertion is at most
$C_kp_0(n^{-c_2}v)^k$ when the old test has raw scale $p_0$.
If a later calculation divides by $v^k$, these factors cancel exactly.
The deterministic window above can therefore be chosen first, then a
fixed depth $k$, and then the finitely many moment orders required by
H\"older.  This statement concerns actual-path marginals; it makes no
claim about redefining conditional means along a one-row interpolation.
\end{remark}

\subsection{Integrated overlap bounds and the conditional error}

The estimates in this subsection are evaluated at a minimizer of
\eqref{eq:6}.  Whenever a small scale $x$ is mentioned, outer regularity
above $x$ is an explicit hypothesis, as are
$x\gtrsim a$ and $e(x)x^2\gtrsim nE^2$.

\begin{proposition}[Local second moments]\label{comp:direct-moments}
On every regular outer dyad of mass $u\asymp v$ beyond a sufficiently
large fixed multiple of $x$, for every sufficiently small fixed
$\epsilon>0$,
\begin{equation}
 \left(\operatorname{avg}_{u\asymp v}D_u^2\right)^{1/2}
 \lesssim_\epsilon n^\epsilon\left[
 (ne(v)v)^{-1/3}+(ne(v)v^2)^{-1/2}
 +(n^{-\delta}E/d^3)^{1/2}+n^{\delta/2}u_*\right].
 \tag{18}\label{eq:18}
\end{equation}
Fixed enlargements of a dyad are permitted in this assertion.

For a small direct test impose additionally $e(x)\gtrsim x^2$, and set
$W=\max(x,\sqrt{e(x)})$.  Then $W\lesssim d$ and
$e(v)\gtrsim W^2$ for $v=O(W)$.  On regular dyads of scale at least
a fixed multiple of $W$, including comparable dyads at $W$ that lie in
the regular region, the right side of \eqref{eq:18} is polynomially
smaller than $x^2/W$.  Between $Cx$ and $O(W)$ it is polynomially
smaller than $x^2/(W^2v)^{1/3}$.  Moreover
\begin{equation}
 \int_0^{CW}D_u\,du=o(x^2),\qquad D=o(x).
 \tag{19}\label{eq:19}
\end{equation}
The integral is truncated at mass $1$, and $C$ may be any fixed
constant.  The second assertion means normalized mean-square
convergence on every compact band
$r\le Cx$, $cx\le u\le Cx$, both in mass and in radius:
\[
 x^{-3}\int_{\{r\le Cx,\,cx\le u\le Cx\}}D^2\,du\longrightarrow0,
 \qquad
 x^{-3}\int_{\{r\le Cx,\,cx\le\alpha(r)\le Cx\}}D^2\,dr
 \longrightarrow0.
\]
If $W=o(1)$, there is a split $*$ in a regular band of scale $W$ with
$S_*\asymp W$, satisfying \eqref{eq:17}, and with $D_*$ polynomially
smaller than $x^2/W$.  It can be chosen above any fixed multiple of $x$.
\end{proposition}

\begin{proposition}[Conditional comparison error]\label{comp:conditional-error}
Let $\lambda\ge n^{-10}$ and
$A=A_\lambda=n^\delta\sqrt n E/\sqrt\lambda$.  If $A=o(1)$,
assume that a minimizer has a regular outer region above $A$.
Then its integral \eqref{eq:7} is at most $n^{10\delta}E$ for all
sufficiently large $n$.  If $A$ stays bounded below by a positive
constant, the same conclusion holds without a regularity hypothesis.

Consequently, if these hypotheses hold along an admissible sweep from
zero to any strength at most $d^2$, then
$0\le\mathcal M_\lambda-f_n(b_0,h_0)\le n^{20\delta}E$.
The fixed-strength assertions apply also to the wall variant, whereas
this last conclusion requires an admissible sweep.
\end{proposition}

We prove both propositions by an averaging estimate that retains cap
components as complete sampling units.  This avoids replacing a
whole-cap trace estimate by a pointwise one.  The two-sided level
averaging argument is a local form of the quantitative Parisi
comparison technique in \cite[Lemma~3.3]{AL2026}.

\begin{lemma}[Trace bounds on sampling units]\label{comp:unit-trace}
Consider a group at comparable masses $u\asymp v$ of total mass at
most $Cv$, and suppose the radii and the raw overlap moments in that
group are bounded at scale $y_0$.  Away from the floor, take as units
complete cap components with comparable endpoint masses, and mass
elements of the contact set.  Let $e_g$ be a lower bound on $e$ in
the group.  If a cap has no initial-wall loss, then within each unit
\begin{equation}
 \operatorname{avg}_{\mathrm{unit}}
       \mathbb E\operatorname{tr}H_u^2
 \lesssim\frac{n^{o(1)}}{ne_g}.
 \tag{20}\label{eq:20}
\end{equation}
On contact this holds pointwise almost everywhere.  In particular the
unit averages of $\mathbb E\operatorname{tr}C_u^2$ and $\mathbb EW_u$
are at most $n^{o(1)}/(ne_gv^2)$ and $n^{o(1)}/(ne_gv)$, respectively.
\end{lemma}

\begin{proof}
In radius coordinates \eqref{eq:1} gives
$dS/dr\ge n \partial_r h\mathbb E\operatorname{tr}H^2$ on continuous parts,
with nonnegative contributions through jumps.  On contact apply
\eqref{eq:9}.  On a cap of radius length $\ell$, integrate and use
$\Delta S\le n^{o(1)}\ell$.  The density $M(u)$ varies by at most a
constant factor within a unit with comparable endpoint masses, and
its mass is comparable to $M(u)\ell$.  Multiplication by this density
and division by unit mass give \eqref{eq:20}.
The inequalities $H\ge uC$ imply
$\operatorname{tr}C^2\le u^{-2}\operatorname{tr}H^2$ and
$\operatorname{tr}HC\le u^{-1}\operatorname{tr}H^2$, proving the last
claims.  These trace comparisons follow directly by taking traces
against positive semidefinite matrices, without a commutativity
assumption.
\end{proof}

\begin{lemma}[Two-sided level averaging]\label{comp:bracket}
In the setting of Lemma~\ref{comp:unit-trace}, let $\ell_0$ be a common
polynomial-scale upper bound for cap widths, including a possible wall
term when it is covered by the trace estimate.  Suppose on the group,
for every fixed $p$, the raw bounds are
$\|X\|_p\le C_p\sqrt{y_0}$ and
$\|R\|_p\le C_py_0$.  Then, up to a factor $n^{O(\epsilon')}$ for
any fixed $\epsilon'>0$, and an edge error $O(k_0n^{-10})$,
\begin{equation}
 \int_{\mathrm{group}}D_u^2\,du
 \lesssim v\left[
       (y_0/k_0)^2+y_0\ell_0+\ell_0^2
              +\frac{k_0+1}{ne_gv^2}\right],
 \qquad 1\le k_0\le n^2.
 \tag{21}\label{eq:21}
\end{equation}
In particular the first and last variable terms may be replaced by
$C[(y_0/(ne_gv^2))^{2/3}+1/(ne_gv^2)]$.
\end{lemma}

\begin{proof}
We give the sampling construction explicitly.  Order the units in
radius, and assign a cap unit of physical mass $\mu$ an effective
interval of length $\mu$.  At every effective position in that interval,
sample its physical point uniformly in the cap's normalized mass.
When taking a side average, use the physical average on that cap,
weighted by the fraction of its effective interval that the side
occupies.  Contact elements retain their physical order.  For a finite-cap approximation, retain each included cap as a whole
sampling unit and keep the contact part as a separate integral; omitted
cap mass tends to zero before taking the limit in $n$, and floor mass is
treated separately.  All auxiliary samples choose locations on the same
planted path and use its same terminal spin.
The effective measure has exactly the original mass.  Physical samples
within a cap may be taken independently, even for partial effective
intervals.  Side averages themselves are deterministic weighted
averages of the random variables $L_u$, not random weights.

Bin units according to their end value of $S$, into at most $k_0$
successive height intervals of width $Cy_0/k_0$.  A cap stays in one
bin; by \eqref{eq:10}, its internal variation of $S$ is at most
$n^{o(1)}\ell_0$.  Within each bin discard effective strips of length $n^{-10}$ at both
ends from the set of evaluation positions only.  For each remaining
effective position $i$, both side averages use the entire portion of the
bin on the corresponding side, including the discarded endpoint strips;
thus both available side lengths are at least $n^{-10}$.
Their deterministic means differ from $\mathbb EL_i=S_i/2$ by at most
$C(y_0/k_0+n^{o(1)}\ell_0)$.

If one side has length $m_1$, its physical weight on a unit of mass
$\mu$ is $q/(m_1\mu)$, where $0\le q\le\mu$ is the effective mass
of that unit included in the side.  By \eqref{eq:1} and the last
assertion of Lemma~\ref{comp:unit-trace}, its variance is at most
\[
 \frac{n^{o(1)}}{ne_gv\,m_1^2}
       \sum_{\mathrm{units}}\frac{q^2}{\mu}
 \le\frac{n^{o(1)}}{ne_gv\,m_1}.
\]
The corresponding integral formula applies on contact.  Integrating
over $i$ costs $O(k_0\log n)$ from the reciprocal side length.

We next compare $L_i$ to its side averages.  For physically ordered
points $j>i$,
\[
 L_j-L_i=(y-X_j)\cdot(X_j-X_i)+\frac12|X_j-X_i|^2.
\]
The first term has second moment at most
\[
 (\mathbb E\operatorname{tr}C_j^2)^{1/2}
       [S_j-S_i+C(D_i+D_j)].
\]
Indeed the conditional covariance at $j$ bounds its square by
$\mathbb E[(X_j-X_i)^{\mathsf T}C_j(X_j-X_i)]$ and then
Cauchy--Schwarz applies.  The identity
\[
 D_u^2=\operatorname{Var}|X_u|^2
       +2\mathbb E X_u^{\mathsf T}C_uX_u
       +\mathbb E\operatorname{tr}C_u^2
\]
bounds the $L^2$ deviations of $|X_u|^2$ and of
$X_i\cdot X_j=\mathbb E_j(X_i\cdot y)$ by the indicated $D$ terms.
This proves the displayed estimate for $\mathbb E|X_j-X_i|^4$ needed
in the Cauchy--Schwarz step.

For the upper deviation of $L_i$, use the right side: the quadratic
term has favorable sign.  For its lower deviation use the left side
and the same identity with the larger physical index last.  Different
units are physically ordered.  If two samples in the same cap have
the reverse physical order, retain also the quadratic term.  Since
$\mathbb E|X_j-X_i|^2=|S_j-S_i|\le n^{o(1)}\ell_0$, interpolation
with an arbitrarily high fixed raw moment gives
\[
 \mathbb E|X_j-X_i|^4\le C_{\epsilon'}n^{\epsilon'}y_0\ell_0.
\]
To see the power accounting, at moment $2q$ the interpolation loss is
at most $(y_0/\ell_0)^{1/(q-1)}$ times logarithms.  Both scales are
bounded by fixed powers of $n$, so a sufficiently large fixed $q$
gives the stated $n^{\epsilon'}$ loss.  A polynomial upper bound
$\ell_0$ may be used even for a much shorter individual cap.

Integrate the squared one-sided estimates and use Jensen for each
side average.  The side kernels have logarithmic marginal-density
loss at either position: in effective coordinates this is the
integral of the reciprocal available side length, cut off at
$n^{-10}$.  A partial unit is spread uniformly across its physical
mass.  Consequently a trace term at either position is bounded by
its whole-unit average in \eqref{eq:20}, with at most this logarithmic
loss.  This remains true when an ordering indicator is present,
because it can be discarded in an upper bound for a nonnegative term.

Put $I_D=\int D^2\,du$ and
$\tau_C=n^{o(1)}/(ne_gv^2)$.  The terms involving one $D$ and one
square root of a trace are bounded by
$n^{o(1)}(I_Dv\tau_C)^{1/2}$.  Young's inequality, with a sufficiently
small inverse-logarithmic coefficient, absorbs the $I_D$ part.  Terms
involving the deterministic height difference are bounded by
$n^{o(1)}v[(y_0/k_0)^2+\ell_0^2+\tau_C]$.
The reverse-order quadratic term adds
$n^{O(\epsilon')}vy_0\ell_0$.  The side-average variances add
$n^{o(1)}k_0/(ne_gv)$.
Finally
\[
 D^2\le4\operatorname{Var}L+\mathbb E\operatorname{tr}C^2,
\]
because $\operatorname{Var}L
=\mathbb E X^{\mathsf T}CX+\tfrac14\operatorname{Var}|X|^2$.
This proves \eqref{eq:21}, including the discarded-strip error by
bounded spin observables.  Optimizing the integer $k_0$ proves the last
assertion: use $k_0$ comparable to $(y_0^2ne_gv^2)^{1/3}$ when this
exceeds one, and otherwise use one.  In the applications $y_0$ may
always be capped by a fixed constant, so this choice is below $n^2$.
\end{proof}

\paragraph{Broad cap units and the initial wall.}
We record the modifications needed outside comparable cap components.
Slice a cap that crosses a large mass ratio into pieces of mass
comparable to $v$, at $u\asymp v$, discarding mass $O(u_*)$ near zero
if necessary.  At any dyad there are only a bounded number of such
broad components: they are ordered disjoint intervals crossing one
of its bounding mass levels.  Integrating the clock inequality over
the whole cap bounds the trace in a slice at the cost
\[
 \frac{M(O(v))}{v}
       \bigl(n^{-\delta}L\log n
                       +r_{\min}\mathbf1_{\mathrm{initial}}\bigr)
 \le n^{o(1)}(1+v/u_*).
\]
For the length term, use
$d^3Lu_*/E\le\sqrt n E/d^2\le1$; for the wall term use
$r_{\min}\le u_*/M(u_0)$ and
$M(v)/M(u_0)\le C[(v+u_*)/(u_0+u_*)]^2$.
Average the observable $L_u$ over the whole slice, without height bins.
Its mean shift is at most $n^{o(1)}L$, with $L$ the deterministic scale
in \eqref{eq:5}, because all points remain in the original cap.
The unordered comparison just proved therefore yields
\eqref{eq:21} without its first term, with $\ell_0=n^{o(1)}L$, and
with the final term replaced by
\[
 \frac{n^{o(1)}(1+v/u_*)}{ne_gv^2}.
\]
For an initial cap with comparable endpoint masses, either its mass
is less than $u_*$, in which case we discard it, or
$M(O(v))r_{\min}\lesssim u_*\lesssim v$, which gives the ordinary
unit bound.  These alternatives cover all wall losses.

\begin{proof}[Proof of \eqref{eq:18}]
In a regular outer band, complete cap units have comparable endpoint
masses: their radius widths are $o(x)$ and both coordinates are
comparable to $v$.  Assign each unit to a dyad with a fixed enlargement.
Use $y_0=Cv$, $e_g\asymp e(v)$, and
$v\ell_0\lesssim n^{-\delta}E/d^3$ in \eqref{eq:21}.
The two optimized variable terms become
$(ne(v)v)^{-2/3}$ and $(ne(v)v^2)^{-1}$ after normalization by the
band mass.  Since $\ell_0=o(v)$, its squared term is absorbed by
$v\ell_0$.  Floor mass in this band is at most $Cv k(cv)$, and raw
moments bound its normalized contribution by $Ck(cv)v^2$.
As $v\gg n^{2\delta}u_*$, its square root is at most
$Cn^{\delta/2}u_*$.  Choose the interpolation loss small relative to
the specified $\epsilon$ and absorb logarithms.  This proves
\eqref{eq:18}.
\end{proof}

\begin{proof}[Proof of Proposition~\ref{comp:conditional-error}]
First treat regular dyads with $v\gtrsim A$ when $A=o(1)$.
Multiply the squared bounds underlying \eqref{eq:18} by $Ce(v)v$,
and include the mismatch using \eqref{eq:9}--\eqref{eq:10}.
The four resulting contributions are bounded, up to arbitrarily small
power losses, by
\[
 n^{-2/3}(e(v)v)^{1/3},\qquad \frac1{nv},\qquad
 n^{-\delta}\frac{e(v)v}{d^3}E,\qquad
 n^\delta e(v)v u_*^2.
\]
Here $e(v)v\le Cd^3\log n$.  Thus the first is at most
$n^{-m+o(1)}E$, the second is at most $E$ because $v\gg u_*$,
the third is at most $n^{-\delta+o(1)}E$, and the last divided by
$E$ is at most
$n^{\delta+o(1)}d^3/(n^2E^3)\lesssim n^{-3m+\delta+o(1)}$.
The extra squared cap-width mismatch is smaller than the cap term.

The remaining region has $r,u=O(A)$, $e\le\lambda$, and
$e_g\gtrsim n^{-\delta}\lambda$.  The last assertion follows from
$\lambda A/d^3\le n^\delta$.  Its raw moment scale can be taken
as $y_0=C\min(A,1)\le CA$.  If $A$ is bounded below, this uses only
the trivial moment bounds and covers the whole interval without an
outer-regularity assumption.  Discard mass $O(u_*)$ near zero and the
floor of total mass $O(u_*)$; their entire error is at most
\[
 C\lambda A^2u_*=Cn^{2\delta}E.
\]
On the remaining comparable mass units apply \eqref{eq:21}, with
$\ell_0\le n^{o(1)}L$.  For $v\gtrsim u_*$ the variable terms,
multiplied by $\lambda v$, are bounded by
\[
 Cn^{4\delta/3+o(1)}E+Cn^{\delta+o(1)}E,
\]
using $\lambda A^2=n^{1+2\delta}E^2$ and
$e_g\gtrsim n^{-\delta}\lambda$.  The mismatch costs are at most
$n^{o(1)}\lambda vAL\le n^{2\delta+o(1)}E$, since $v\lesssim A$
and $L\le u_*$.  Broad units use the preceding paragraph; their
trace term costs at most
$n^{\delta+o(1)}(1+v/u_*)/(nv)\lesssim n^{\delta+o(1)}E$.
Initial-wall terms were either discarded at the same cost or included
in that trace estimate and in \eqref{eq:10}.

There are $O(\log n)$ dyads.  Cap units crossing the division between
inner and outer regions can be included by a fixed enlargement of the
inner region; all estimates above tolerate that change.  Choose
$\epsilon'$ in \eqref{eq:21} sufficiently small compared with
$\delta$, and absorb the logarithmic losses.  The whole integral
\eqref{eq:7} is at most $n^{10\delta}E$.
Lemma~\ref{comp:envelope} then permits integration over an interval
of logarithmic strength of length $O(\log n)$, with the negligible
initial interval treated directly.  This proves the $n^{20\delta}E$
assertion under its explicitly stated sweep hypotheses.
\end{proof}

\begin{proof}[Proof of the direct-test conclusions of
Proposition~\ref{comp:direct-moments}]
The additional condition $e(x)\gtrsim x^2$ implies $x\lesssim d$ and
$W\lesssim d$.  Since $\lambda\le d^2$, the argument of $f$ at a
radius of order $W$ is bounded, and $e(v)\gtrsim W^2$ there.  Above
$W$, $e(v)v\gtrsim W^3$.  The first two terms in \eqref{eq:18},
divided by $x^2/W$ on such dyads, are bounded by
$Cn^{-1/3}/x^2$ and $Cn^{-1/2}/x^3$; both have a polynomial saving
since $x\gtrsim n^{-1/6+m}$.  The cap term divided by $a^2/d$ is
bounded by
\[
 C\sqrt{n^{-\delta}\frac{d^3}{n^2E^3}},
\]
and the floor term has ratio at most
$Cn^{\delta/2}d^3/(n^2E^3)$.  These also have polynomial savings.
Since $x^2/W\gtrsim a^2/d$, they prove the claimed high-band bound.
Between $Cx$ and $O(W)$ use $e(v)\gtrsim W^2$.  Dividing the first
two terms by $x^2/(W^2v)^{1/3}$ gives respectively
$Cn^{-1/3}/x^2$ and at most $Cn^{-1/2}/x^3$; the remaining terms are
handled by the same comparison with $a^2/d$.  Take $\epsilon$ smaller
than these margins.  The savings are uniform over the dyads of a
sequence with fixed comparability constants.

We next prove the mass version of the second assertion in
\eqref{eq:19}.  On $cx\le u\le Cx$, $r\le Cx$, use the ordinary
or broad-unit estimates with $y_0=Cx$, $v\asymp x$, and
$e_g\asymp e(x)$.  Units extending towards zero can always be sliced
as above.  Since
\[
 ne_gx^4\gtrsim(x/u_*)^2\longrightarrow\infty
\]
polynomially, both optimized variable terms in \eqref{eq:21} are
$o(x^2)$ after normalization by mass.  The broad trace loss is also
harmless:
\[
 \frac{1+x/u_*}{ne_gx^4}
 \lesssim (u_*/x)^2+u_*/x=o(1).
\]
The cap and wall scales satisfy
$\ell_0/x+r_{\min}/x\le n^{-\delta+o(1)}L/x=o(1)$.
The floor's raw contribution is $o(x^3)$ because its mass is
$O(u_*)=o(x)$.  This proves the asserted mass mean square.

For the radius conclusion a direct argument is needed; one cannot
discard a radius interval because it carries little physical mass.
Let
$I_0=\{r\le Cx:cx\le\alpha(r)\le Cx\}$, an interval up to endpoints.
Raw moments give $S,D\le C_1x$ and $\|L\|_p\le C_px$ there.
The total increase of $S$ is $O(x)$, so \eqref{eq:1}, with
$\partial_r h\ge e_g/O(\log n)$, gives
\[
 \int_{I_0}\mathbb E\operatorname{tr}H^2\,dr
       \le\frac{n^{o(1)}x}{ne_g},\qquad
 x^{-1}\int_{I_0}\mathbb E\operatorname{tr}C^2\,dr
       \le\frac{n^{o(1)}}{ne_gx^2}.
\]
For any radius subinterval $I\subset I_0$, apply the last line of
\eqref{eq:1} with deterministic mass weight
$\mathbf1_I(r(u))/(x\alpha'(r(u)))$.  It is legitimate at fixed $n$
because the floor is strictly positive.  Since
$\mathbb EW\le\mathbb E\operatorname{tr}H^2/(cx)$,
\[
 \operatorname{Var}\left(x^{-1}\int_I L\,dr\right)
 \le Cx^{-2}\int_I\frac{\mathbb EW}{\alpha'}\,dr
 \le\frac{n^{o(1)}}{ne_gx^2k(Cx)}=o(x^2).
\]
Only one reciprocal density is paid: one power cancels against
$du=\alpha' dr$.  More explicitly,
$n^{2\delta}u_*\ll x$ and
$k(Cx)\gtrsim n^\delta u_*^2/x^2$, so the last variance divided by
$x^2$ is at most
\[
 Cn^{1-\delta+o(1)}\frac{E^2}{e_gx^2}
 \le n^{-\delta+o(1)}.
\]
This remains true on a subinterval that lies entirely on the floor.
Also $x^{-3}\int\mathbb E\operatorname{tr}C^2\,dr=o(1)$ by the
preceding trace bound.

Partition $I_0$ into at most $k_0$ occupied intervals by deterministic
$S$ heights, each of span $O(x/k_0)$.  Remove strips of radius length
$\tau x$ at both ends of each interval.  Their contribution to
$x^{-3}\int D^2\,dr$ is $O(k_0\tau)$ by raw moments.  Every remaining
point has left and right side averages of radius length at least
$\tau x$.  For fixed $k_0,\tau$ their variances are uniformly $o(x^2)$
by the direct radius bound, and their means differ from $\mathbb EL_i$
by $O(x/k_0)$.  The ordered identity used in
Lemma~\ref{comp:bracket} has the favorable quadratic sign on the two
sides.  Its remaining first term has square at most
$Cx(\mathbb E\operatorname{tr}C_j^2)^{1/2}$, or the analogous
quantity at $i$, using the raw fourth moment of the projection
increment.  The side kernels have size at most $1/(\tau x)$.
Their integrated error is $o(x^3)$ by Cauchy--Schwarz and the
integrated trace bound.  Consequently
\[
 \limsup_{n\to\infty}x^{-3}\int_{I_0}D^2\,dr
 \le C(k_0^{-2}+k_0\tau).
\]
Here $k_0,\tau$ are fixed before the limit in $n$.  Choose $k_0$
arbitrarily large, then $\tau$ sufficiently small depending on it.
This proves the radius part of \eqref{eq:19}, including virtual
splits on floor intervals.

To prove the first assertion of \eqref{eq:19}, the dyads between
$Cx$ and $CW$ contribute at most
\[
 o(x^2)W^{-2/3}\int_x^{CW}v^{-1/3}\,dv=o(x^2).
\]
The savings used here are polynomial and uniform, so logarithmic
dyadic losses are harmless.  On $[\varepsilon x,Cx]$, for fixed
$\varepsilon>0$, the mass mean-square assertion and Cauchy--Schwarz
give $o(x^2)$.  On $[0,\varepsilon x]$ the raw overlap bound gives
at most $C\varepsilon x^2$.  First let $n\to\infty$ and then
$\varepsilon\downarrow0$.  Mass omitted from the optimization costs
only $O(xu_0)=o(x^2)$; the raw bound there was proved by projection
from the regular region.

Finally suppose $W=o(1)$.  In a regular band at a sufficiently large
fixed multiple of $W$, \eqref{eq:18} and the mismatch estimates give
$|S-r|=o(W)$ in normalized mass measure off the floor; cap widths are
$o(W)$ and floor mass is negligible.  Monotonicity then supplies a
subinterval of mass comparable to $W$ on which $S\asymp W$.
Apply Lemma~\ref{comp:good-split}.  The averaged $D$ bound on this
band has a polynomial saving relative to $x^2/W$.  Choose $c_1$
smaller than the corresponding polynomial measure-saving exponent.
Chebyshev's inequality then makes the set failing a still
polynomially saved $D$ bound smaller than the central good window.
Select a point outside that set.  The all-moment estimate holds at
every deterministic point in the window, so this extra selection
preserves \eqref{eq:17} without choosing a moment-dependent window.
Increasing the fixed multiple of $W$ places the split above any
prescribed fixed multiple of $x$.
\end{proof}

\begin{remark}[Logical use of the estimates]\label{comp:dependency}
Propositions~\ref{comp:direct-moments} and
\ref{comp:conditional-error} have used only the path identities,
the constrained first variation, and their stated outer-region
hypothesis.  In particular, no global comparison error has been
assumed to obtain \eqref{eq:19} or the good split.  The one-site
calculation uses these estimates to identify the possible direct-test
profiles.  The finite propagation argument then establishes outer
regularity throughout \eqref{eq:8}; only at that point is the sweep
conclusion of Proposition~\ref{comp:conditional-error} invoked without
an unproved regularity premise.
\end{remark}
